% docs/manual/user/appendices/B-cli-reference.tex
% 附录 B：CLI 命令完整参考

\chapter{CLI 命令完整参考}

\openbench{} CLI 入口为 \cli{openbench}（安装后自动注册到 PATH）。共 7 个主命令 + 1 个 \cli{version}。本附录给出每个命令的签名、参数、退出码与示例。

\section{全局选项}

\begin{minted}{text}
openbench [OPTIONS] COMMAND [ARGS]...

Options:
  --version  Show version and exit
  --help     Show help and exit
\end{minted}

CLI 使用 lazy command loading；启动 \cli{openbench --help} 不会加载子命令模块。

\section{\cli{openbench version}}

\textbf{Synopsis}:

\begin{minted}{bash}
openbench version
\end{minted}

\textbf{用途}：打印版本号。

\textbf{退出码}：始终 0。

\textbf{示例输出}：

\begin{minted}{text}
openbench 3.0.0a1
\end{minted}

\section{\cli{openbench init}}

\textbf{Synopsis}:

\begin{minted}{bash}
openbench init [-o OUTPUT]
\end{minted}

\textbf{选项}：

\begin{itemize}
  \item \cli{-o, --output PATH}：输出 yaml 路径（默认 \file{openbench.yaml}）
\end{itemize}

\textbf{用途}：交互式向导生成最小 \file{openbench.yaml}。详见第 2 章。

\textbf{退出码}：

\begin{itemize}
  \item 0：成功生成
  \item 非 0：用户中断或写入失败
\end{itemize}

\section{\cli{openbench check}}

\textbf{Synopsis}:

\begin{minted}{bash}
openbench check CONFIG
\end{minted}

\textbf{参数}：

\begin{itemize}
  \item \texttt{CONFIG}：\file{.yaml} / \file{.yml} 配置路径（必须存在）
\end{itemize}

\textbf{用途}：零数据访问的预飞行校验。详见第 6 章。

\textbf{退出码}：

\begin{itemize}
  \item 0：所有校验通过
  \item 1：YAML 语法、schema、reference 解析任一失败
\end{itemize}

\section{\cli{openbench run}}

\textbf{Synopsis}:

\begin{minted}{bash}
openbench run CONFIG [OPTIONS]
\end{minted}

\textbf{参数}：

\begin{itemize}
  \item \texttt{CONFIG}：YAML 配置路径
\end{itemize}

\textbf{选项}：

\begin{itemize}
  \item \cli{--dry-run}：只校验 + 列任务，不实际计算
  \item \cli{--cores N}：覆盖 \yamlkey{project.num_cores}（必须 ≥ 1）
  \item \cli{--variables VAR}：可重复，仅评估指定变量子集
  \item \cli{--remote HOST}：[当前未实现] 远程执行
  \item \cli{--dump-config}：把中间 namelist 与 runner 配置写到 \file{output/<run>/debug/}
  \item \cli{--comparison-only}：跳过 evaluation，仅基于已有结果重做 comparison
\end{itemize}

\textbf{退出码}：

\begin{itemize}
  \item 0：评估全部成功
  \item 1：配置加载失败 / \cli{--remote} 触发的 NotImplemented / 评估全失败
  \item 部分失败时根据 \texttt{results.status} 决定（partial 也返 1）
\end{itemize}

详见第 6 章。

\section{\cli{openbench data}}

数据集管理命令组。子命令：

\subsection{\cli{openbench data list}}

\begin{minted}{bash}
openbench data list [--variable VAR]
\end{minted}

列出全部 reference 数据集；\cli{--variable} 仅列含该变量的。详见第 4 章。

\subsection{\cli{openbench data download}}

\begin{minted}{bash}
openbench data download NAMES...
\end{minted}

\warninline{当前版本（v3.0a1）该子命令是 stub —— 打印 "Dataset download not yet implemented" 后退出。计划在后续版本实现。}

\subsection{\cli{openbench data status}}

\begin{minted}{bash}
openbench data status
\end{minted}

报告 registry 中数据集数量。

\warninline{Local cache 检测部分尚未实现。}

\subsection{\cli{openbench data path}}

\begin{minted}{bash}
openbench data path NAME
\end{minted}

打印某数据集的本地路径。支持 base 名展开多分辨率变体。详见第 4 章。

\textbf{退出码}：1 当 \texttt{NAME} 在 catalog 中不存在。

\subsection{\cli{openbench data register}}

\begin{minted}{bash}
openbench data register NAME --root-dir PATH \
  [--data-type {grid|stn}] [--tim-res {Month|Day|Hour|Year}] \
  [--grid-res FLOAT] [--category CATEGORY] \
  [--years START END] [--fulllist CSV_PATH] \
  [-v "Var:nc_name:unit"]... \
  [-f "Var:fb_name:fb_unit:conversion"]...
\end{minted}

向 user-level catalog 注册新 reference 数据集。详见第 4 章。

\textbf{交互模式}：不带 \cli{-v} 时进入交互式变量录入。

\section{\cli{openbench model}}

模型 profile 管理命令组。

\subsection{\cli{openbench model list}}

\begin{minted}{bash}
openbench model list
\end{minted}

列出全部 model profile（21 个内置）。详见第 5 章。

\subsection{\cli{openbench model show}}

\begin{minted}{bash}
openbench model show NAME
\end{minted}

显示 profile 的完整变量映射（含 NC 名、单位、fallback、compute 表达式）。

\textbf{退出码}：1 当 profile 不存在。

\subsection{\cli{openbench model register}}

\begin{minted}{bash}
openbench model register NAME [OPTIONS]
\end{minted}

\textbf{选项}：

\begin{itemize}
  \item \cli{--data-type \{grid|stn\}}
  \item \cli{--grid-res FLOAT}
  \item \cli{--tim-res \{Month|Day|Hour|Year\}}
  \item \cli{--description TEXT}
  \item \cli{-v "VarName:nc\_name:unit"}（可重复）
  \item \cli{-f "VarName:fb\_name:fb\_unit:conversion"}（可重复）
  \item \cli{--append-only}：仅添加新变量，不覆盖已有
\end{itemize}

详见第 5 章。

\subsection{\cli{openbench model remove-var}}

\begin{minted}{bash}
openbench model remove-var MODEL_NAME VAR_NAME
\end{minted}

从 profile 中删除一个变量。

\section{\cli{openbench migrate}}

\textbf{Synopsis}:

\begin{minted}{bash}
openbench migrate OLD_CONFIG [-o OUTPUT]
\end{minted}

\textbf{参数}：

\begin{itemize}
  \item \texttt{OLD\_CONFIG}：v1.x \file{.nml} / v2.0 \file{.json} / 早期 \file{.yaml} 路径
\end{itemize}

\textbf{选项}：

\begin{itemize}
  \item \cli{-o, --output}：新 yaml 路径（默认 \file{openbench.yaml}）
\end{itemize}

详见第 9 章。

\section{\cli{openbench gui}}

\textbf{Synopsis}:

\begin{minted}{bash}
openbench gui [CONFIG] [--remote]
\end{minted}

\textbf{参数}：

\begin{itemize}
  \item \texttt{CONFIG}（可选）：启动时加载已有配置
\end{itemize}

\textbf{选项}：

\begin{itemize}
  \item \cli{--remote}：以远程模式启动（要求 \cli{[remote]} extra）
\end{itemize}

\textbf{依赖}：必须装 \cli{[gui]} extra（PySide6）。详见第 8 章。

\textbf{退出码}：

\begin{itemize}
  \item 0：GUI 正常退出
  \item 1：依赖检查失败（未装 PySide6）
\end{itemize}

\section{脚本化使用要点}

\begin{itemize}
  \item 所有命令都遵循 Unix 退出码约定，可与 shell 流控结合
  \item 日志默认到 stderr；解析时用 \texttt{2>\&1} 合并或 \texttt{2> log.err} 分离
  \item \cli{run} 是阻塞命令；后台运行用 \texttt{\&} 或 nohup
  \item 不支持配置文件之外的 stdin 输入；必要时先 \cli{init} 写文件再 \cli{run}
\end{itemize}
